\documentclass[11pt,a4paper]{article}
\usepackage[margin=24mm]{geometry}
\usepackage[T1]{fontenc}
\usepackage{lmodern}
\usepackage{amsmath,amssymb}
\usepackage{booktabs,tabularx,array}
\usepackage{graphicx}
\usepackage{tikz}
\usetikzlibrary{arrows.meta,positioning,calc}
\usepackage{pgfplots}
\usepgfplotslibrary{groupplots}
\pgfplotsset{compat=1.18}
\usepackage[font=small,labelfont=bf]{caption}
\usepackage{placeins}
\usepackage{microtype}
\usepackage{titlesec}
\titlespacing*{\section}{0pt}{2.5ex plus 1ex}{1.3ex}
\titlespacing*{\subsection}{0pt}{2ex plus .5ex}{1ex}
\usepackage[round,authoryear]{natbib}
\usepackage[hidelinks,breaklinks]{hyperref}
\usepackage{xurl}
\definecolor{steel}{HTML}{315D7B}
\definecolor{rust}{HTML}{AE5939}
\definecolor{cluster0}{HTML}{347CA8}
\definecolor{cluster1}{HTML}{BF6D20}
\definecolor{cluster2}{HTML}{378951}
\definecolor{cluster3}{HTML}{BE4545}
\definecolor{cluster4}{HTML}{855AA5}
\definecolor{cluster5}{HTML}{8B6657}
\definecolor{cluster6}{HTML}{B85D8B}
\definecolor{cluster7}{HTML}{64757B}
\definecolor{cluster8}{HTML}{A39037}
\definecolor{cluster9}{HTML}{168A94}
\tikzset{box/.style={draw=steel!70,fill=steel!5,rounded corners=1pt,align=center,text width=3.1cm,minimum height=1.2cm,font=\small,inner sep=6pt},widebox/.style={box,text width=12.4cm,minimum height=1cm},arrow/.style={-{Latex[length=2mm]},thick,draw=steel!80}}
\setlength{\parindent}{0pt}
\setlength{\parskip}{6pt}
\setlength{\emergencystretch}{2em}
\renewcommand{\topfraction}{.95}
\renewcommand{\bottomfraction}{.9}
\renewcommand{\textfraction}{.05}
\renewcommand{\floatpagefraction}{.75}
\newcommand{\diagram}[1]{\resizebox{\linewidth}{!}{\input{figures/#1.tex}}}

\begin{document}
\section*{Three cross-sectional study diagram atlas}
Native TikZ and PGFPlots sources for the revised manuscript. All numerical figures use the saved clean-run tables; schematic diagrams are identified in their captions.
\clearpage
\begin{figure}[p]
\centering
\diagram{population}
\caption{Index populations and continuity. The flow population has 39 complete histories; every date-specific fit retains all 41 instruments.}\end{figure}
\clearpage
\begin{figure}[p]
\centering
\diagram{windows}
\caption{Dependence across the three cross-sections. Bars depict trailing annual market windows schematically; fundamentals follow their recorded issuer reporting dates, not these uniform windows.}\end{figure}
\clearpage
\begin{figure}[p]
\centering
\diagram{pipeline}
\caption{Repeated cross-sectional study workflow. The preprocessing reference carries forward, while each date receives a fresh representation and partition.}\end{figure}
\clearpage
\begin{figure}[p]
\centering
\diagram{alignment}
\caption{Matching clusters with an observed overlap example. Five of the eight earlier group-1 members occupy the later group-1 match. This local example illustrates the similarity calculation; the actual assignment optimises all group pairs jointly.}\end{figure}
\clearpage
\begin{figure}[p]
\centering
\diagram{quality}
\caption{The same partitions evaluated in two geometries. UMAP has the stronger embedding silhouette at every date; this advantage does not consistently hold in the underlying financial features.}\end{figure}
\clearpage
\begin{figure}[p]
\centering
\diagram{seed_variability}
\caption{Ten pairwise adjusted Rand indices from five UMAP seeds at each date. Points are offset horizontally for readability; red segments mark the means. Pairs share fitted partitions and are dependent, so the spread is not a confidence interval or an empirical null for temporal switches.}\end{figure}
\clearpage
\begin{figure}[p]
\centering
\diagram{flows}
\caption{Native LaTeX rendering of all 39 complete instrument histories: 19 stable and 20 shifted. Bold tickers change aligned label at least once. APN, MNP, GLN and PAN are excluded only from this complete-history view. Labels identify aligned groups, not industries; vertical positions are layout.}\end{figure}
\clearpage
\begin{figure}[p]
\centering
\diagram{jaccard}
\caption{Jaccard overlap among the 40 continuing members in each adjacent comparison. Each cell uses the common eligible universe for that period. Numeric values are given explicitly; a zero denotes no shared continuing member.}\end{figure}
\clearpage
\begin{figure}[p]
\centering
\diagram{profile0}
\caption{June 2025 aligned-cluster median deviations from the same-date median, in June 2025 MinMax units. Counts include all 41 date-specific instruments. Colours use the same range across the three figures; gaps in cluster numbers reflect retired identifiers.}\end{figure}
\clearpage
\begin{figure}[p]
\centering
\diagram{profile1}
\caption{December 2025 aligned-cluster median deviations from the same-date median, in June 2025 MinMax units. Counts include all 41 date-specific instruments. Colours use the same range across the three figures; gaps in cluster numbers reflect retired identifiers.}\end{figure}
\clearpage
\begin{figure}[p]
\centering
\diagram{profile2}
\caption{June 2026 aligned-cluster median deviations from the same-date median, in June 2025 MinMax units. Counts include all 41 date-specific instruments. Colours use the same range across the three figures; gaps in cluster numbers reflect retired identifiers.}\end{figure}
\clearpage
\begin{figure}[p]
\centering
\diagram{feature_sensitivity}
\caption{Complete-history outcomes under four feature representations using the UMAP recipe. Each bar contains the same 39 continuously observed instruments; four partial histories are outside every bar. The counts are representation-dependent and are not estimates of true economic change.}\end{figure}
\clearpage
\begin{figure}[p]
\centering
\diagram{cluster_count}
\caption{Cluster-count sensitivity redrawn from the archived numerical table for the same clean run. Blue evaluates financial features and red evaluates the embedding. Eight groups remain the historical specification, not a validated natural count.}\end{figure}
\clearpage
\begin{figure}[p]
\centering
\diagram{evidence}
\caption{Interpretation of separate evidence streams. Group movement, feature profiles and distance flags can motivate follow-up, but are not combined into an automatically confirmed anomaly. Stochastic, measurement and surrogate limitations apply to the interpretation.}\end{figure}
\end{document}
